{\Large\textbf{1\quad Hidden Charge}}

{\large\textbf{1.1\quad Introduction}}

An unknown point charge $Q$ is fixed in a region of space. Electrons launched
parallel to the $z$ axis far from the charge will scatter electrostatically off
of the fixed charge and strike a detecting screen. It is possible learn about
the details of the hidden charge by varying the initial kinetic energy as well
as the initial $x_\mathrm{i}$ and $y_\mathrm{i}$ coordinates of the electron
beam and measuring the final coordinates $x_\mathrm{f}$ and $y_\mathrm{f}$ of
where an electron strikes a finite flat screen perpendicular to the $z$ axis and
located at $z = 0$.

It is useful to know the Rutherford scattering formula,
\[
b = \frac{kqQ}{2E}\frac{1}{\tan(\theta/2)}
\]
where $b$ is the impact parameter, $E$ is the energy of the electron,
$q = -1.602 \times 10^{-19}$C is the charge of the electron,
$k = 8.99 \times 10^{9}$ Nm$^2$/C$^2$, and $\theta$ is the scattering angle.
The impact parameter is defined as the closest approach of the electron to the
target, assuming that the electron were unaffected by the target and hence would
move in a straight line; the scattering angle is angle between the original
velocity vector of the electron far from the target and the final velocity
vector of the electron far from the target after scattering.

\begin{center}\includegraphics[width=0.50\linewidth]{figures/EuPhO_2020_Q4_fig1.png}\end{center}

{\large\textbf{1.2\quad Task}}

The task is to determine the position $(x_Q, y_Q, z_Q)$ and also the magnitude
and sign of the fixed charge $Q$, as precisely as possible. You should provide
rough, order of magnitude error estimates on these results. There is Gaussian
error associated with initial beam location that is on the order of $0.5$ mm.

As with all experiments, you must provide clearly labelled tables of data,
clearly labelled graphs, and sufficient formulae derivations to make it clear
what you have measured, and how you are deriving your results.

{\large\textbf{1.3\quad Program Interface}}

The program asks for an accelerating voltage with the prompt

\texttt{Beam accelerating voltage in V:}

Enter a number between 1 and 10000, and press \textbf{return}. The program then
asks for the initial launch coordinates, starting with $x_\mathrm{i}$, with the
prompt

\texttt{x-coordinate of the electron beam in cm:}

Enter a number between -20 and 20 and then press \textbf{return}. Finally, the
program asks for $y_\mathrm{i}$, with the prompt

\texttt{y-coordinate of the electron beam in cm:}

Enter a number between -20 and 20 and then press \textbf{return}. If you enter
an invalid number for any of these three, the program will prompt you with

\texttt{Invalid entry.}

and will then prompt you for the value again, reminding you of the allowed
limits.

After the three numbers have been entered, the program will output

\texttt{Electron beam fired with parameters (x, y, V) =}

and it will restate your entered values, and then

\texttt{Electron detected at (x, y) =}

and give the screen location of the detected electron. However, if the electron
misses the finite size screen, you will be told

\texttt{Electron not detected...}

The program then repeats, allowing you to enter in a new set of initial
coordinates.
